\section{混合分布モデル}
\subsection{授業内容}
\subsubsection{科目の中でのこのコマの位置づけ}
ここまでGLMM，HLMとさまざまな分布を組み合わせて利用するモデルについてみてきた。

今回は混合正規分布モデルと0過剰ポアソンモデルを考える。これまではデータが全体的に均質的であることを想定していたが，そもそも異なる種類のデータが混じり合っていると考えられる場合，
\subsubsection{コマ主題細目}
\begin{description}
\end{description}
\subsubsection{キーワード}
\begin{itemize}
\end{itemize}
\subsection{授業情報}
\paragraph{コマの展開方法}
講義/遠隔可/演習
\subsubsection{予習・復習課題}
\paragraph{予習}反復測定デザインの分析例を確認しておく。そこで分布が混合されていることを確認しておく。
\paragraph{復習}階層線形モデルが応用できるようなデータ例を身の回りから考えてみるとよい。その上で，データがどのような背景から生成されているかの設計図を書き，設計図からコードに起こすという手順を一歩ずつ確認しておこう。